% !TEX root = ../Projektdokumentation.tex
\section{Projektplanung} 
\label{sec:Projektplanung}

\subsection{Projektphasen}
\label{sec:Projektphasen}
Für die Umsetzung des Projektes standen den Autoren 70 Schulstunden zur Verfügung. Diese wurden vor Projektbeginn auf
verschiedene Phasen verteilt, die während des Entwicklungsprozesses durchlaufen werden. Da das Projektthema aus Lernfeld
10 (Blockwoche 1) fortgeführt wird, befindet sich das Projekt zu diesem Zeitpunkt bereits im Entwicklungsprozess. Aus der
Tabelle~\bhref{tab:Zeitplanung} kann deshalb eine Zeitplanung für die noch zu erledigenden Aufgaben ab Beginn dieses
Projektes entnommen werden.
\tabelle{Zeitplanung}{tab:Zeitplanung}{Zeitplanung} \\

\subsection{Ressourcenplanung}
\label{sec:Ressourcenplanung}
Eine detaillierte Übersicht der verwendeten Ressourcen lässt sich im Anhang
\bhref{app:VerwendeteRessourcen} nachlesen. Dort sind alle Ressourcen aufgelistet, welche für das 
Projekt eingesetzt wurden. Das beinhaltet Hard- und Softwareressourcen sowie Personal. Bei der
Auswahl der verwendeten Software wurde darauf geachtet, dass diese kostenfrei zur Verfügung steht
oder uns bereits Lizenzen bereit gestellt waren. 

\subsection{Entwicklungsprozess}
\label{sec:Entwicklungsprozess}
Für die Umsetzung des Projekts wurde ein an das Wasserfallmodell angelehnter Entwicklungsprozess gewählt. Dieses Vorgehensmodell zeichnet
sich durch eine klare, sequentielle Struktur der einzelnen Projektphasen aus und eignet sich insbesondere für Projekte mit weitgehend stabilen
Anforderungen.

Zu Beginn des Projekts erfolgte die Analyse der Ausgangssituation sowie die Definition der Anforderungen. Auf dieser Grundlage wurde ein
Soll-Konzept erarbeitet, das als Basis für die anschließende Umsetzung diente.

Im weiteren Verlauf wurde die Implementierung der Anwendung durchgeführt. Dabei wurden die zuvor definierten Anforderungen schrittweise umgesetzt.
Aufgrund des schulischen Projektumfelds erfolgte die Entwicklung überwiegend ohne formale Iterationszyklen, wobei kleinere Anpassungen und Korrekturen
während der Implementierungsphase berücksichtigt wurden.

Nach der Implementierung wurden grundlegende Funktionstests durchgeführt, um die korrekte Funktionsweise der Anwendung sicherzustellen. Hierbei lag
der Fokus insbesondere auf der Überprüfung zentraler Prozesse wie der Bestellung, der Statusverwaltung und der Bestandsaktualisierung.

Das gewählte Vorgehensmodell ermöglichte eine strukturierte Bearbeitung des Projekts und eine klare Trennung der einzelnen Projektphasen.
Gleichzeitig zeigte sich, dass bei zukünftigen Projekten mit dynamischeren Anforderungen auch iterative Vorgehensmodelle, wie beispielsweise agile
Methoden, in Betracht gezogen werden könnten.